\section{Discussion, limitations and scope}\label{sec:disc}
\textbf{What the evidence supports.} In our simulations a linear regression of outcomes on own treatment and treated-neighbour fraction is the most reliable regression GATE estimator across mechanisms: its noise error is exactly $\sigma\|w\|_2$ and scales as $\sqrt{d/n}$ (R2); its bias is governed by how curved the exposure--response is \emph{near the exposure values the design generates} (R3), so it is small under cluster designs with good clusters and under linear-in-exposure interference, and large under Bernoulli designs with curved responses. Replacing the linear model by flexible learners or by many features did not help (R5); neither did the correct model form when the design leaves the nonlinear region unexcited (R4). These results suggest that the practical lever is the \emph{design} (clusters, treatment probability, saturation) together with realised-assignment diagnostics $\|w\|_1,\|w\|_2$, rather than model sophistication.

\textbf{Practical guidance (conditional on our simulations).} (i) Compute $\|w\|_2$ and $\|w\|_1$ from the realised assignment; $\hat\sigma\|w\|_2$ is a reliable noise-error scale, and a large $\|w\|_1$ with a curved response flags extrapolation risk. (ii) At high degree relative to $n$ ($d/n$ large) and no good clusters, regression may be worse than DM; interference needs to be a substantial share of the GATE (here $\gtrsim25\%$) for regression to pay off. (iii) Use cluster randomisation only if the graph has a small edge cut for the clusters; otherwise its variance cost is not repaid. (iv) Use cluster-robust SEs on the randomisation clusters rather than i.i.d./HC SEs, and treat intervals as optimistic when bias is plausible. (v) Be sceptical of the graph: corrupted edges bias Bernoulli-design estimates roughly proportionally to the corruption.

\textbf{Limitations.} (1) \emph{Synthetic outcomes}: the mechanisms (eight, all of ``one-hop function of exposure'' type except \texttt{twohop}/\texttt{peer}) are our own choices, and all conclusions are conditional on them; real interference may differ. Absolute error levels depend on our baseline heterogeneity and noise scale (the assignment-induced variance from baselines dominated the noise variance), and on the chosen normalisation ($\mathrm{GATE}=2$). (2) \emph{One population per configuration}: graphs, baselines and community effects are drawn once (seed 0) and the replicates vary only assignment and noise; variability across graph draws is not quantified, and the 6 graphs are few. The reported CIs thus reflect Monte-Carlo, not population, uncertainty. (3) The error analysis (Prop.~\ref{prop:decomp}) is exact for OLS, but Prop.~\ref{prop:bern} is an asymptotic statement for regular graphs with independent neighbour treatments; we verified it numerically with a finite-$d$ binomial version (all 54 comparisons within 2.5 s.e., two beyond 1.96; at $\bar d\ge64$ the MC error is large), and it does not cover heavy-tailed degrees (BA, Stanford3), where the relevant $d$ is not obvious. (4) \emph{Estimators not studied}: we did not implement the full methods of \citet{chin2019} (resampling inference) or \citet{han2023} (covariate selection and block bootstrap), graph neural networks \citep{kipf2017}, staggered rollouts \citep{cortez2022}, or covariate-adjusted/regularised variants that could reduce the variance floor or the cluster-level baseline variance; GBDT was used with one fixed hyper-parameter setting. The H\'ajek baseline uses a single exposure threshold ($q{=}0.75$) and thus targets a different estimand. (5) We use one clustering algorithm (Louvain at a hand-picked resolution per graph, chosen once before the main runs; Exp.~10 studies the resolution) and equal-probability cluster-level assignment. (6) Exp.~6 and 3 use different SBM block counts ($K{=}15$ and $30$); comparisons across them should not be made. (7) Replicate counts (100--500) limit the resolution of small biases ($\lesssim0.03$). (8) Compute: all experiments used 4 CPU cores; no hyper-parameter tuning was done on test outcomes.

\textbf{Negative and unexpected results.} (a) Regression adds little over DM under cluster randomisation with good clusters (and DM is best on ego-Facebook), so the accuracy gain is mostly due to the design. (b) Flexible learners were worse than the simple linear model. (c) The oracle model was far from the best in some conditions (R4). (d) Standard SEs are badly miscalibrated. (e) $\|w\|_1$ is a poor general predictor of bias (R9).
